\section{Robustez e diagnóstico dos modelos}
\label{sec:robustez}

Duas baterias de verificação fecham a fase de estimação.

\textbf{Sem o biênio pandêmico.} Todos os modelos foram reestimados
excluindo 2020 e 2021, os anos em que a casuística e os denominadores
sofreram os choques documentados na primeira
parte. A Figura~\ref{fig:robustezpandemia} compara os contrastes OSS vs
Direta nas duas amostras: as diferenças são mínimas (o tempo de
permanência passa de 10,2\% para 11,4\% de queda, a produção de 43,7\%
para 44,7\% de alta, e assim por diante). \textbf{Nenhum dos
contrastes principais depende dos anos de pandemia.}

\begin{table}[H]
\centering
\caption{Contraste OSS \emph{versus} Direta estimado no painel completo
e na amostra que exclui 2020 e 2021, nos sete modelos principais. Cada
bloco traz o \textbf{efeito} já traduzido para a unidade de leitura, ao
lado do valor bruto do modelo e da sua margem de erro. A leitura
pretendida é a comparação entre os dois blocos: quanto mais parecidos,
menos o resultado depende dos anos de pandemia.}
\label{tab:sempandemia}
\small
\setlength{\tabcolsep}{4pt}
\begin{tabular}{lrrrrrr}
\toprule
& \multicolumn{3}{c}{Painel completo} & \multicolumn{3}{c}{Sem 2020 e 2021} \\
\cmidrule(lr){2-4}\cmidrule(lr){5-7}
Modelo & Efeito & Coef. & EP & Efeito & Coef. & EP \\
\midrule
Mortalidade (geral)        & 5,8\% $\to$ 4,6\% & $-0{,}2237$ & 0,1311 & 5,6\% $\to$ 4,4\% & $-0{,}2329$ & 0,1157 \\
Fração alta complexidade   & 8,9\% $\to$ 7,6\% & $-0{,}2044$ & 0,2348 & 8,9\% $\to$ 7,7\% & $-0{,}1937$ & 0,2333 \\
Faturamento real por saída & $-8{,}3\%$  & $-0{,}0866$ & 0,0948 & $-8{,}0\%$  & $-0{,}0830$ & 0,0978 \\
Tempo de permanência       & $-10{,}2\%$ & $-0{,}1080$ & 0,0468 & $-11{,}4\%$ & $-0{,}1207$ & 0,0527 \\
Produção (saídas)          & $+43{,}7\%$ & $0{,}3629$  & 0,0895 & $+44{,}7\%$ & $0{,}3694$  & 0,0892 \\
Ocupação de internação     & $+22{,}7\%$ & $0{,}2044$  & 0,0779 & $+22{,}3\%$ & $0{,}2015$  & 0,0753 \\
Ocupação de UTI            & $+19{,}3\%$ & $0{,}1769$  & 0,0603 & $+22{,}1\%$ & $0{,}1995$  & 0,0600 \\
\quad {\footnotesize cálculo alternativo} & $+22{,}8\%$ & $0{,}2058$ & 0,0747 & $+24{,}3\%$ & $0{,}2175$ & 0,0723 \\
\bottomrule
\end{tabular}
\end{table}

\comofeito{cada linha é um dos sete modelos principais, e cada bloco de
três colunas é uma amostra: à esquerda o painel inteiro, à direita a mesma
estimação depois de retirar 2020 e 2021. \textbf{Efeito} é o resultado
traduzido para a unidade que o gestor usa. Nos indicadores de quantidade
ele é a variação percentual; nos dois que já são taxas, a mortalidade e a
fração de alta complexidade, ele é o par de taxas que o deslocamento liga,
e a seta se lê ``sai de'' e ``chega a'', com a taxa de partida sendo a
média observada da administração Direta naquela amostra. \textbf{Coef.} é o
mesmo resultado no valor bruto do modelo, que não se lê diretamente, e é
por isso que a coluna Efeito existe. \textbf{EP} é a margem de erro,
calculada agrupando as observações por hospital, porque onze anos do mesmo
hospital não são onze informações independentes. Nenhum efeito muda de
sinal e nenhum muda de magnitude de forma relevante, que é o que a figura
seguinte mostra em forma gráfica. A ocupação de UTI aparece em
\textbf{duas linhas} porque foi calculada de duas maneiras: a primeira é a
adotada no documento, escolhida depois que os testes de aderência do
Apêndice~\ref{ap:diagnosticos} mostraram que ela descreve melhor estes
dados; a segunda usa a mesma conta das outras seis linhas e fica na tabela
para que a coluna possa ser lida de cima a baixo na mesma régua.}
\newpage
\comoler{cada linha traz duas estimativas do mesmo efeito: o círculo
azul escuro usa todos os anos e o quadrado verde exclui 2020 e 2021.
Quanto mais próximos os dois símbolos, mais robusto o resultado; aqui
eles praticamente se sobrepõem em todas as linhas.}

\begin{figure}[H]
\centering
\includegraphics[width=.82\textwidth]{fig_est_05_robustez_sem_pandemia.png}
\caption{Contraste OSS vs Direta estimado no painel completo e na
amostra sem 2020 e 2021, para os modelos em escala percentual.}
\label{fig:robustezpandemia}
\end{figure}

\comofeito{gráfico de pontos. Cada linha é um indicador (eixo vertical) e
o eixo horizontal traz o \textbf{efeito OSS vs Direta} em variação
percentual. Cada linha traz \textbf{dois símbolos do mesmo efeito}: um
estimado no painel completo e outro na amostra que exclui 2020 e 2021.
Quanto mais próximos os dois símbolos, mais o resultado independe dos
anos de pandemia.}

\textbf{Aderência dos modelos.} A segunda bateria confere se cada modelo
de fato descreve os dados que ele resume. Todos passam, e o material
completo, com o diagnóstico de cada um e os gráficos correspondentes,
está no Apêndice~\ref{ap:diagnosticos}.

% A terceira bateria, "Descasamento de calendário nos denominadores", saiu
% daqui em 04/08/2026, a pedido expresso de João, junto com a sua tabela
% (a antiga Tabela 10, `tab:calendario`) e as duas ressalvas que a fechavam.
%
% O QUE ELA MEDIA, para quem precisar reabrir o assunto: as duas ocupações
% dividem produção acumulada no ano por capacidade cadastrada vezes 365, e
% quando numerador e denominador não cobrem o mesmo período a razão mede
% calendário e não desempenho. A bateria não corrigia o problema — ele não é
% corrigível na granularidade anual —, ela media o TETO dele, reestimando os
% contrastes sem os conjuntos que continham as observações afetadas. O
% resultado era tranquilizador em toda a escada: excluir os 30 anos parciais
% de UTI move o contraste de 19,3% para 19,8% (0,06 erro-padrão), e no degrau
% mais severo os dois indicadores de ocupação se deslocam em DIREÇÕES
% OPOSTAS (+0,31 e −0,38 erro-padrão), assinatura de ruído de amostragem e
% não de viés sistemático, que empurraria ambos para o mesmo lado.
%
% Foi por esse resultado que a bateria saiu: ela custava três páginas e uma
% tabela de dez linhas para concluir que nada do que o documento afirma
% depende do assunto. Toda ela continua reproduzível por
% `analises/verificar_bateria_calendario.py` e `auditar_uti_ano_parcial.py`,
% e o registro da decisão está na seção 21 de PENDENCIAS.md.
Fecham o quadro as ressalvas permanentes de leitura. A identificação
do efeito de gestão dentro de hospital repousa em \textbf{cinco
conversores}, três deles com um único ano após a mudança; as margens
de erro agrupadas com tão poucas unidades tratadas tendem a ser
otimistas, e os efeitos devem ser lidos como ordens de grandeza. A
categoria Privado (3 CNES) permaneceu como controle de absorção em
todos os modelos, sem interpretação. E os resultados de mortalidade e
ocupação do estudo de eventos carregam a contaminação de tendências
anteriores já descrita, que o tempo de permanência não apresenta.
